% from ega1-5-fr.tex
\oldpage[I]{137}
prouve (ii). Vu (i) et (ii), (iii) et (iv) sont équivalentes
(\hyperref[I.3.5.1-fr]{3.5.1}), et il suffit de démontrer (iv). Or,
$X_{(S')}\times_{Y_{(S')}}X_{(S')}$ s'identifie canoniquement à
$(X\times_Y X)\times_Y Y_{(S')}$ en vertu de
(\hyperref[I.3.3.11-fr]{3.3.11}) et de
(\hyperref[I.3.3.9.1-fr]{3.3.9.1}), et on vérifie aussitôt que le morphisme
diagonal $\Delta_{X_{(S')}}$ s'identifie alors à
$\Delta_X\times_Y 1_{Y_{(S')}}$~; la proposition résulte donc de
(\hyperref[I.4.3.1-fr]{4.3.1}).

Pour établir (v), considérons, comme dans
(\hyperref[I.5.3.13-fr]{5.3.13}) la factorisation
$X\xrightarrow{\Gamma_f}X\times_Z Y\xrightarrow{p_2}Y$ de $f$, en remarquant
que $p_2=(g\circ f)\times_Z 1_Y$~; l'hypothèse que $g\circ f$ est séparé
entraîne que $p_2$ est séparé par (iii) et (i), et comme $\Gamma_f$ est une
immersion, $\Gamma_f$ est séparé par (i), donc $f$ est séparé par (ii).
Enfin, pour démontrer (vi), rappelons que les préschémas
$X_{\mathrm{red}}\times_{Y_{\mathrm{red}}}X_{\mathrm{red}}$ et
$X_{\mathrm{red}}\times_Y X_{\mathrm{red}}$ s'identifient canoniquement
(\hyperref[I.5.1.7-fr]{5.1.7})~; si on désigne par $j$ l'injection
$X_{\mathrm{red}}\to X$, le diagramme
\[
  \xymatrix{
    X_{\mathrm{red}}\ar[r]^{\Delta_{X_{\mathrm{red}}}}\ar[d]_j &
    X_{\mathrm{red}}\times_Y X_{\mathrm{red}}\ar[d]^{j\times_Y j}\\
    X\ar[r]_{\Delta_X} & X\times_Y X
  }
\]
est commutatif (\hyperref[I.5.3.15-fr]{5.3.15}), et la proposition résulte de
ce que les flèches verticales sont des homéomorphismes des espaces
sous-jacents (\hyperref[I.4.3.1-fr]{4.3.1}).

\begin{corollary}[5.5.2]
\label{I.5.5.2-fr}
Si $f:X\to Y$ est séparé, la restriction de $f$ à tout sous-préschéma de $X$
est séparée.
\end{corollary}

Cela résulte de (\hyperref[I.5.5.1-fr]{5.5.1, (i) et (ii)}).

\begin{corollary}[5.5.3]
\label{I.5.5.3-fr}
Si $X$, $Y$ sont deux $S$-préschémas tels que $Y$ soit séparé au-dessus de
$S$, $X\times_S Y$ est séparé au-dessus de $X$.
\end{corollary}

C'est un cas particulier de (\hyperref[I.5.5.1-fr]{5.5.1, (iv)}).

\begin{proposition}[5.5.4]
\label{I.5.5.4-fr}
Soit $X$ un préschéma, et supposons que son espace sous-jacent soit réunion
d'une famille finie de parties fermées $X_k$ $(1\leq k\leq n)$~; on considère
pour chaque $k$ le sous-préschéma réduit de $X$ ayant $X_k$ pour espace
sous-jacent (\hyperref[I.5.2.1-fr]{5.2.1}) et on le note encore $X_k$. Soit
$f:X\to Y$ un morphisme, et pour chaque $k$, soit $Y_k$ une partie fermée de
$Y$ telle que $f(X_k)\subset Y_k$~; on note encore $Y_k$ le sous-préschéma
réduit de $Y$ ayant $Y_k$ pour espace sous-jacent, de sorte que la restriction
$X_k\to Y$ de $f$ à $X_k$ se factorise en
$X_k\xrightarrow{f_k}Y_k\to Y$ (\hyperref[I.5.2.2-fr]{5.2.2}). Pour que $f$
soit séparé, il faut et il suffit que les $f_k$ le soient.
\end{proposition}

La nécessité résulte de
(\hyperref[I.5.5.1-fr]{5.5.1, (i), (ii) et (v)}). Inversement, si la condition
de l'énoncé est satisfaite, chacune des restrictions $X_k\to Y$ de $f$ est
séparée (\hyperref[I.5.5.1-fr]{5.5.1, (i) et (ii)})~; si $p_1$, $p_2$ sont les
projections de $X\times_Y X$, le sous-espace $\Delta_{X_k}(X_k)$ s'identifie
au sous-espace $\Delta_X(X)\cap p_1^{-1}(X_k)$ de l'espace sous-jacent à
$X\times_Y X$ (\hyperref[I.5.3.16-fr]{5.3.16})~; ces sous-espaces étant
fermés dans $X\times_Y X$, il en est de même de leur réunion $\Delta_X(X)$.

Supposons en particulier que les $X_k$ soient les
\emph{composantes irréductibles} de $X$~; alors on peut supposer que les $Y_k$
sont des composantes irréductibles de $Y$
(\hyperref[0.2.1.5-fr]{0, 2.1.5})~; la prop.
(\hyperref[I.5.5.4-fr]{5.5.4}) ramène donc dans ce cas la notion de séparation
au cas des préschémas \emph{intègres} (\hyperref[I.2.1.7-fr]{2.1.7}).

